\documentclass[pre,preprint,preprintnumbers,aps,eqsecnum,superscriptaddress,notitlepage,showpacs,floatfix]{revtex4-1}

\usepackage{graphicx}
\usepackage{dcolumn}
\usepackage{bm}
\usepackage{amsmath}
\usepackage{amssymb}
\usepackage{verbatim}
\usepackage{color}
\usepackage{hyperref}
\usepackage{tikz}
\usepackage{bbold}

\newcommand{\dd}{\,d}
\newcommand{\Tr}{\operatorname{Tr}}
\newcommand{\avg}[1]{\left\langle #1 \right\rangle}
\newcommand{\del}{\delta}

\begin{document}

\author{ccc}
\address{Mathematical Biology Section, LBM, NIDDK, NIH}
\title{The Effective Action Approach to Neural Networks with Random Connectivity}
\date{\today}

\begin{abstract}
We formulate an effective action and two-particle-irreducible (2PI)
framework for neural networks with random connectivity.  The construction is
organized around the microscopic action and the vertices generated by
connectivity disorder, nonlinear gain, phase-density transport, reset or
boundary conditions, and discrete spike events.  This organization makes clear
that a 2PI closure is not unique: different approximations correspond to
different truncations of the diagrammatic expansion of the 2PI effective
action.  The usual Sompolinsky--Crisanti--Sommers closure is recovered by
retaining only the disorder bubble and a smooth Gaussian output covariance.
For spiking networks, the current scalar approximation used in the code treats
the filtered spike-count shot noise as a homogeneous contribution that inflates
the equal-time drive variance.  We also distinguish the smooth-feedback,
shot-renormalized, and operational finite-branch notions of criticality.
\end{abstract}

\maketitle

%=====================================================================
\section{Introduction}
%=====================================================================

Randomly connected neural networks provide a useful setting in which to
compare dynamic mean-field theory, path-integral methods, and systematic
effective-action approximations.  In rate networks, the large-$N$ limit leads
to the Sompolinsky--Crisanti--Sommers (SCS) self-consistent equation for the
single-neuron autocorrelation.  In spiking networks the output of a neuron is
a flux of discrete events through threshold, possibly accompanied by reset,
refractory dynamics, or phase wrapping.  The path integral therefore contains
more vertices than the rate-network MSR action alone.

The numerical code in this repository keeps several closures available.  The
current default scalar phase closure is the \emph{inflated-initial-condition}
closure.  This document uses the same conventions as the code modules
\texttt{scripts/rn\_rate.py}, \texttt{scripts/rn\_phase.py}, and
\texttt{scripts/rn\_binary.py}.

%=====================================================================
\section{Rate-network action and SCS limit}
\label{sec:rate_action}
%=====================================================================

Consider the random rate network
\begin{align}
\dot u_i(t)+u_i(t)-\sum_{j=1}^N w_{ij}f(u_j(t))=0,
\qquad
w_{ij}\sim N\!\left(0,\frac{\sigma^2}{N}\right).
\label{eq:rate_network}
\end{align}
The MSR representation of the deterministic dynamics is
\begin{align}
S_{\rm MSR}
=
\sum_i\int\dd t\;\tilde u_i(t)
\left[
\dot u_i(t)+u_i(t)-\sum_j w_{ij}f(u_j(t))
\right].
\label{eq:rate_MSR_before_avg}
\end{align}
Averaging over Gaussian random weights gives
\begin{align}
S_{\rm rate}
=
\sum_i\int\dd t\;\tilde u_i(t)[\dot u_i(t)+u_i(t)]
-
\frac{\sigma^2}{2}\sum_i\int\dd t\dd s\;
\tilde u_i(t)\tilde u_i(s)R(t,s),
\label{eq:rate_action_avg}
\end{align}
where
\begin{align}
R(t,s)=\frac1N\sum_j f(u_j(t))f(u_j(s)).
\label{eq:R_rate_def}
\end{align}
In this convention $R$ excludes the factor of $\sigma^2$.  The effective input
covariance is $Q(t,s)=\sigma^2R(t,s)$.

The rate-network SCS closure treats the single-neuron input as a Gaussian
process with covariance $C_{11}(t,s)=\avg{u(t)u(s)}$ and evaluates
\begin{align}
R(t,s)=\avg{f(u(t))f(u(s))}_{\rm Gaussian}.
\label{eq:R_gaussian_rate}
\end{align}
For a stationary solution with $\tau=t-s$ and unit decay rate,
\begin{align}
\frac{d^2}{d\tau^2}C_{11}(\tau)
=
C_{11}(\tau)-\sigma^2 R\bigl(C_{11}(\tau);C_0\bigr),
\qquad C_0=C_{11}(0).
\label{eq:SCS_unit_beta}
\end{align}
With synaptic decay rate $\beta$,
\begin{align}
\frac{d^2}{d\tau^2}C_{11}(\tau)
=
\beta^2\left[C_{11}(\tau)-\sigma^2 R\bigl(C_{11}(\tau);C_0\bigr)\right].
\label{eq:SCS_beta}
\end{align}
For the monotone rate branch, the energy condition is
\begin{align}
H(C_0)
=
C_0^2-2\sigma^2\int_0^{C_0}R(x;C_0)\,\dd x=0.
\label{eq:rate_energy}
\end{align}

%=====================================================================
\section{General 2PI effective action}
\label{sec:general_2PI}
%=====================================================================

Let the collective field be
\begin{align}
\xi_\alpha\in\{u,\tilde u,\eta,\tilde\eta,\ldots\}.
\end{align}
Expanding the action around $a_\alpha=\avg{\xi_\alpha}$ gives
\begin{align}
S[\xi]
=S[a]
+\frac12 S^{(2)}_{\alpha\beta}[a]\delta\xi_\alpha\delta\xi_\beta
+\frac{1}{3!}S^{(3)}_{\alpha\beta\gamma}[a]
\delta\xi_\alpha\delta\xi_\beta\delta\xi_\gamma
+\cdots .
\label{eq:vertex_expansion}
\end{align}
The 2PI effective action is
\begin{align}
\Gamma[a,C]
=
S[a]
+\frac12\Tr\ln C^{-1}
+\frac12\Tr\left(S^{(2)}[a]C\right)
+\Gamma_2[a,C]
+{\rm const}.
\label{eq:2PI_general}
\end{align}
The term $\Gamma_2[a,C]$ is the sum of closed two-particle-irreducible vacuum
diagrams built from interaction vertices and the full propagator $C$.  The
stationarity equations imply the Dyson equation
\begin{align}
C^{-1}=L-\Sigma,
\qquad
L_{\alpha\beta}=S^{(2)}_{\alpha\beta}[a],
\qquad
\Sigma_{\alpha\beta}=2\frac{\delta\Gamma_2}{\delta C_{\alpha\beta}}.
\label{eq:dyson_general}
\end{align}
A closure is therefore specified by a chosen truncation of $\Gamma_2$.

%=====================================================================
\section{Spiking action and conventions}
\label{sec:spiking_action}
%=====================================================================

Consider phase or voltage-like variables
\begin{align}
\dot v_i &= F_i(u_i,v_i),\\
\dot u_i &= -\beta u_i
+
\beta\sum_j\left(w_{ij}-\lambda\bar w\right)\nu_j(t),
\label{eq:u_dynamics_spiking}
\end{align}
where $\lambda\in\{0,1\}$ selects row-sum correction and
$\nu_j(t)=\sum_l\delta(t-t_j^l)$ is the spike train.  The empirical density is
\begin{align}
\eta_i(v,t)=\delta(v_i(t)-v).
\end{align}
Away from threshold and reset boundaries,
\begin{align}
\partial_t\eta_i+\partial_v(F_i\eta_i)=0.
\label{eq:klimontovich}
\end{align}
The threshold flux is
\begin{align}
\nu_i(t)=F_i(u_i,v_T)\eta_i(v_T,t).
\end{align}

The spiking action is organized as
\begin{align}
S=S_\eta+S_u+S_{\rm reset}+S_{\rm spike}+S_{\rm dis}.
\label{eq:full_spiking_action_decomp}
\end{align}
The transport part is
\begin{align}
S_\eta
=
\sum_i\int\dd t\dd v\;
\tilde\eta_i
\left[\partial_t\eta_i+
\partial_v(F_i\eta_i)-\delta(t-t_0)\rho_i(v)\right].
\label{eq:Seta_transport}
\end{align}
The synaptic filtering term before disorder averaging is
\begin{align}
S_u
=
\sum_i\int\dd t\;\tilde u_i[\dot u_i+\beta u_i]
-
\beta\sum_{ij}\int\dd t\;\tilde u_i
\left(w_{ij}-\lambda\bar w\right)\nu_j(t).
\label{eq:Su_before_avg_clean}
\end{align}
After averaging the weights,
\begin{align}
S_{\rm dis}
=
-
\frac{\beta^2\sigma^2}{2}
\sum_i\int\dd t\dd s\;
\tilde u_i(t)\tilde u_i(s)R_\nu(t,s),
\label{eq:Sdis_clean}
\end{align}
with
\begin{align}
R_\nu(t,s)
=
\frac1N\sum_j\nu_j(t)\nu_j(s)
-
\lambda\bar\nu(t)\bar\nu(s).
\label{eq:Rnu_def_clean}
\end{align}
Again $R_\nu$ excludes $\sigma^2$.

For stochastic spike emission or escape, the Doi-Peliti event vertex is
\begin{align}
S_{\rm spike}
=
-
\sum_i\int\dd t\;\lambda_i(t)
\left[
\exp\left(\Delta\tilde\eta_i(t)\right)-1-\Delta\tilde\eta_i(t)
\right],
\label{eq:Sspike_clean}
\end{align}
where $\lambda_i(t)=F_i(u_i,v_T)\eta_i(v_T,t)$ and
$\Delta\tilde\eta_i$ is the response-field increment for the event.  For an
explicit reset $v_T\to v_R$,
\begin{align}
\Delta\tilde\eta_i(t)=\tilde\eta_i(v_R,t)-\tilde\eta_i(v_T,t).
\end{align}
The quadratic term
\begin{align}
S_{\rm spike}^{(2)}
=
-
\frac12\sum_i\int\dd t\;\lambda_i(t)
\left[\Delta\tilde\eta_i(t)\right]^2
\label{eq:Sspike_quadratic}
\end{align}
generates the local event-counting covariance
\begin{align}
\avg{\nu_i(t)\nu_i(s)}_{\rm sing}
=
\avg{\nu_i(t)}\delta(t-s).
\label{eq:spike_train_sing_cov}
\end{align}

For the periodic phase model used in the simulations, phase wrap identifies
$v_T=\pi$ with $v_R=-\pi$.  The scalar closure used in the code does not
retain the full reset-response structure; it approximates the recurrent
feedback through a smooth Gaussian output covariance plus the filtered
shot-noise variance described below.

%=====================================================================
\section{Closure hierarchy used in the code}
\label{sec:closure_hierarchy}
%=====================================================================

The code keeps several closures available, but only the current closure is
used by default.  The relevant diagram classes are:
\begin{center}
\renewcommand{\arraystretch}{1.35}
\begin{tabular}{lll}
\hline
Diagram class & Origin & Physical effect \\
\hline
Disorder bubble & Weight averaging & SCS-type recurrent input \\
Gain vertices & Expansion of $f(u)$ or $F(u,v)$ & Nonlinear covariance feedback \\
Phase-density vertices & $\tilde\eta\partial_v(F\eta)$ & Phase response and resonant kernels \\
Reset/boundary vertices & Threshold-reset or phase wrapping & Conservation through spikes \\
Event-counting vertices & Discrete spike emission & Shot noise and short-time cusp \\
Higher cumulant diagrams & Cubic and quartic vertices & Non-Gaussian corrections \\
\hline
\end{tabular}
\end{center}

\subsection{Smooth Gaussian scalar closure}

The smooth scalar approximation replaces the output by
\begin{align}
g(u)=a_3(v_T)F(u,v_T),
\end{align}
and evaluates the centered covariance
\begin{align}
R_{\rm smooth}(c,C_{\rm tot})
=
\avg{g(u_1)g(u_2)}_{\Sigma(c,C_{\rm tot})}
-
\avg{g(u)}_{C_{\rm tot}}^2,
\label{eq:R_smooth}
\end{align}
where $c$ is the lag covariance and $C_{\rm tot}=C_{11}(0)$ is the fixed
single-time variance.  The corresponding input covariance is
\begin{align}
Q_{\rm smooth}(c,C_{\rm tot})=\sigma^2 R_{\rm smooth}(c,C_{\rm tot}).
\end{align}
The code evaluates this either by direct Gauss-Hermite or QMC quadrature, or
by a stable Hermite expansion.  The Hermite method avoids catastrophic
cancellation in
$\avg{g(u_1)g(u_2)}-\avg{g}^2$ when the variance is large.

\subsection{Current scalar phase closure: inflated initial condition}

The filtered shot-noise covariance generated by a white spike-counting
singularity is
\begin{align}
C_{\rm shot}(\tau)
=
\frac{\beta D_{\rm shot}}{2}e^{-\beta|\tau|},
\qquad
D_{\rm shot}
=
\sigma^2\avg{g(u)}_{C_{\rm tot}}.
\label{eq:Cshot_filtered}
\end{align}
This contribution satisfies the homogeneous filter equation away from
$\tau=0$.  The current scalar closure therefore treats shot noise as an
inflation of the equal-time variance.  The total covariance satisfies
\begin{align}
C_{11}''(\tau)
=
\beta^2\left[
C_{11}(\tau)-Q_{\rm smooth}\bigl(C_{11}(\tau),C_{\rm tot}\bigr)
\right],
\qquad
C_{11}'(0)=0,
\label{eq:inflated_ic_ode}
\end{align}
with
\begin{align}
C_{\rm tot}
=
C_0+\frac{\beta D_{\rm shot}(C_{\rm tot})}{2}.
\label{eq:inflated_ic}
\end{align}
The smooth chaotic amplitude $C_0$ is determined by
\begin{align}
H(C_0;C_{\rm tot})
=
C_0^2-2\sigma^2\int_0^{C_0}R_{\rm smooth}(x,C_{\rm tot})\,\dd x=0.
\label{eq:inflated_energy}
\end{align}
Equations \eqref{eq:inflated_ic}--\eqref{eq:inflated_energy} are solved
self-consistently.  In \texttt{scripts/rn\_phase.py} this is
\texttt{solver="inflated\_ic"}, and it is the default theory in the plotting
scripts.

\subsection{Exploratory generalized kernel}

The scalar phase closure can be supplemented by a low-frequency local kernel,
\begin{align}
\mathcal L_C
=
\partial_\tau^2+2\Gamma\partial_\tau+\Omega_0^2-\beta^2,
\label{eq:local_kernel}
\end{align}
so that
\begin{align}
\mathcal L_CC_{11}(\tau)
=
-\beta^2Q_{\rm smooth}\bigl(C_{11}(\tau),C_{\rm tot}\bigr).
\label{eq:local_kernel_equation}
\end{align}
In the code, \texttt{kernel\_omega} and \texttt{kernel\_damping} are
dimensionless multiples of $\beta$ by default.  This is an exploratory
effective-kernel approximation, not a replacement for the full phase-density
2PI closure.

\subsection{Phase-density advection and threshold-return ringing}

For the phase model, $\partial_v F=0$ does not imply that the density
correlator is trivial.  The phase-density correlator advects around the
circle.  For fixed $u$ the Green's function has the schematic form
\begin{align}
C_{33}(v,v',\tau\mid u)
\sim
\sum_{m\in\mathbb Z}A_m(u)\,
\delta\!\left(v-v'-F(u)\tau+2\pi m\right),
\label{eq:C33_ring_green}
\end{align}
or, equivalently,
\begin{align}
(\partial_\tau+F(u)\partial_v)C_{33}(v,v',\tau\mid u)
=
S_{33}(v,v';u)\delta(\tau).
\label{eq:C33_advection}
\end{align}
The source $S_{33}$ includes the equal-time density covariance.  When the
observable is the threshold flux, every return of the advected density packet
to $v_T=\pi$ injects a singular contribution into the threshold-threshold
correlator.  Thus, for a nearly fixed drive,
\begin{align}
C_{33}(v_T,v_T,\tau\mid u)
\sim
\sum_{m\geq 0}B_m(u)\,
\delta\!\left(\tau-mT(u)\right),
\qquad
T(u)=\frac{2\pi}{F(u)}.
\label{eq:C33_threshold_comb}
\end{align}
This is the microscopic source of the ``ringing'': the phase density crosses
threshold, wraps around the circle, and crosses again, injecting a new
delta-like contribution at each passage through $\pi$.

The spike-output covariance inherits these returns,
\begin{align}
Q_{\rm ring}(\tau\mid u)
\sim
\sigma^2F(u)^2C_{33}(v_T,v_T,\tau\mid u),
\label{eq:Q_ring}
\end{align}
in addition to the local event-counting shot noise at $\tau=0$.  The scalar
Gaussian closure used in the current code replaces this threshold-return
comb by a smooth covariance of $g(u)=\rho F(u)$.  That averaging is most
defensible in a chaotic regime where $u(t)$ fluctuates enough that the return
times $T(u)$ are broadly distributed and the peaks smear together.  Near
threshold, or whenever $u$ is narrow enough that many neurons share a similar
period, the omitted $C_{33}$ ringing is precisely the structure expected to
produce oscillatory wobbles in $C_{11}$.

%=====================================================================
\section{Phase model and criticality}
\label{sec:phase_model}
%=====================================================================

The phase model in the simulations is
\begin{align}
F(u)=\alpha[I+u]_+^{1/\alpha},
\qquad
v_T=\pi,
\qquad
\rho=\frac1{2\pi},
\qquad
g(u)=\rho F(u).
\label{eq:phase_model_F}
\end{align}
With row-sum correction, the mean drive is zero and the relevant smooth
covariance is centered:
\begin{align}
R_{\rm smooth}(c,C_{\rm tot})
=
\avg{g(u_1)g(u_2)}_{\Sigma(c,C_{\rm tot})}
-
\avg{g(u)}_{C_{\rm tot}}^2.
\end{align}
The shot-noise strength is
\begin{align}
D_{\rm shot}(C_{\rm tot})
=
\sigma^2\rho\avg{F(u)}_{C_{\rm tot}}.
\end{align}

The code tracks three different threshold notions:
\begin{enumerate}
\item \emph{Smooth-feedback threshold}
\begin{align}
\sigma_c^{\rm smooth}
=
\frac{1}{\rho F'(0)}
=
\frac{2\pi}{I^{1/\alpha-1}}.
\label{eq:sigma_smooth}
\end{align}
This is the old zero-variance linearization.

\item \emph{Shot-renormalized static threshold}.  This accounts for the
shot-noise-dressed variance below transition:
\begin{align}
1
&=
\sigma_c^2\rho^2
\avg{F'(u)}_{C_*}^2,
\label{eq:sigma_shot_a}\\
C_*
&=
\frac{\beta\sigma_c^2\rho}{2}
\avg{F(u)}_{C_*}.
\label{eq:sigma_shot_b}
\end{align}
These equations are solved by
\texttt{phase\_shot\_renormalized\_criticality}.  The code uses adaptive
one-dimensional quadrature by default because clipping makes
$\avg{F'}$ discontinuous or singular for several $\alpha$ values.

\item \emph{Operational finite-branch threshold}.  This is the largest
$\sigma$ for which the selected nonlinear closure returns a finite stationary
branch.  It is computed by \texttt{phase\_operational\_criticality} and
depends on the chosen approximation.
\end{enumerate}

For $I=\beta=1$, the shot-renormalized static threshold shifts the smooth
threshold substantially.  For example, the current quadrature gives
\begin{center}
\begin{tabular}{c|ccc}
$\alpha$ & $\sigma_c^{\rm smooth}$ & $\sigma_c^{\rm shot}$ &
$\sigma_c^{\rm shot}/\sigma_c^{\rm smooth}$\\
\hline
$0.5$ & $6.28$ & $4.74$ & $0.75$\\
$0.75$ & $6.28$ & $7.57$ & $1.21$\\
$1.0$ & $6.28$ & $10.77$ & $1.71$\\
$1.5$ & $6.28$ & $17.82$ & $2.84$\\
$2.0$ & $6.28$ & $24.86$ & $3.96$
\end{tabular}
\end{center}
These are static criteria, not a proof of the full nonlinear branch structure.

%=====================================================================
\section{Binary neuron network}
\label{sec:binary}
%=====================================================================

For comparison, consider binary neurons with states $n_i\in\{0,1\}$,
\begin{align}
0\xrightarrow{f(u_i)}1,
\qquad
1\xrightarrow{\mu}0,
\qquad
\dot u_i=-\beta u_i+\beta\sum_jw_{ij}n_j.
\end{align}
For a linear gain $f(u)=f_0+f_1u$, define
\begin{align}
\gamma=\mu+f_0,
\qquad
\bar n=\frac{f_0}{\gamma},
\qquad
c_1=\frac{f_1\mu}{\gamma},
\qquad
D_0=2\bar n(1-\bar n)\gamma.
\end{align}
The exact linear 2PI result is
\begin{align}
\hat C_{nn}(\omega)
=
\frac{D_0(\beta^2+\omega^2)}
{(\gamma^2+\omega^2)(\beta^2+\omega^2)-c_1^2\beta^2\sigma^2}.
\label{eq:Cnn_exact}
\end{align}
The denominator factors as
\begin{align}
(\omega^2+\kappa_+^2)(\omega^2+\kappa_-^2),
\qquad
\kappa_\pm^2
=
\frac12\left[(\gamma^2+\beta^2)
\pm
\sqrt{(\gamma^2-\beta^2)^2+4c_1^2\beta^2\sigma^2}
\right].
\end{align}
For $g=c_1\sigma/\gamma<1$,
\begin{align}
C_{nn}(\tau)
=
A_+e^{-\kappa_+|\tau|}+A_-e^{-\kappa_-|\tau|},
\qquad
A_\pm
=
\frac{D_0}{2}
\frac{\beta^2-\kappa_\pm^2}
{\kappa_\pm(\kappa_\mp^2-\kappa_\pm^2)}.
\end{align}
The linear binary transition occurs at $g=1$.

For clipped binary gains, the code keeps both an effective-linear
quasi-static approximation and a direct Gaussian integral closure.  The
integral closure can use a Hermite covariance expansion and optionally adds
the intrinsic telegraph covariance of binary switching.

%=====================================================================
\section{Comparison of active closures}
\label{sec:comparison}
%=====================================================================

\begin{center}
\renewcommand{\arraystretch}{1.4}
\begin{tabular}{lccc}
\hline
 & Rate/SCS & Phase scalar default & Binary linear \\
\hline
Output model
& Smooth gain
& Smooth Gaussian $g(u)$ plus shot-inflated $C(0)$
& Binary jump process \\
Shot noise
& None
& $C_{\rm shot}(0)=\beta D_{\rm shot}/2$
& $D_0$ in numerator \\
Default code path
& \texttt{rn\_rate.theory\_rate\_autocorr}
& \texttt{solver="inflated\_ic"}
& \texttt{theory\_binary\_autocorr}\\
Criticality
& SCS energy condition
& smooth, shot-renormalized, and branch thresholds
& $c_1\sigma/\gamma=1$\\
\hline
\end{tabular}
\end{center}

The central distinction is that shot noise and oscillations arise from
different parts of the action.  Shot noise is produced by event discreteness
and inflates the equal-time drive variance in the current scalar closure.
The microscopic oscillatory mechanism is threshold-return ringing in
$C_{33}(v_T,v_T,\tau)$: each recirculation of phase density through $\pi$
injects a delta-like contribution to the spike-output covariance.  An
effective phase-density kernel is a coarse representation of that omitted
structure.

%=====================================================================
\section{Discussion}
\label{sec:discussion}
%=====================================================================

The effective-action formulation provides a way to organize approximations to
random neural networks without committing prematurely to a unique closure.  In
rate networks, the natural low-order truncation gives the SCS equation.  In
spiking networks, additional vertices from threshold flux, reset or phase
wrapping, and event counting generate shot noise, memory kernels, and possible
oscillatory collective modes.

The current code reflects this hierarchy.  The active scalar phase theory
uses an inflated initial condition, stable Hermite or quadrature evaluation of the
nonlinear gain covariance, and separate criticality diagnostics.  The next
systematic improvement is to build a phase-density closure that retains the
$C_{33}$ threshold-return delta comb, or a controlled effective kernel derived
from it, rather than collapsing those peaks into the scalar Gaussian
covariance.

\begin{thebibliography}{99}

\bibitem{Sompolinsky1988}
H.~Sompolinsky, A.~Crisanti, H.~J.~Sommers,
Phys.\ Rev.\ Lett.\ \textbf{61}, 259 (1988).

\bibitem{Crisanti1987}
A.~Crisanti and H.~Sompolinsky,
Phys.\ Rev.\ A \textbf{36}, 4922 (1987).

\bibitem{BuiceChow2013}
M.~A.~Buice and C.~C.~Chow,
PLOS Comput.\ Biol.\ \textbf{9}, e1002872 (2013).

\end{thebibliography}

\end{document}
